\section{Menus et Navigation (Menu)}

\subsection{Structures \texttt{menu\_item\_config} et \texttt{menubar\_config}}
\begin{lstlisting}[language=C]
typedef struct menu_item_config {
    const char *label;
    menu_item_kind kind;
    const struct menu_item_config *children;
    size_t child_count;
    bool enabled;
    menu_item_callback on_activate;
    gpointer user_data;
} menu_item_config;

typedef struct {
    const menu_section_config *sections;
    size_t section_count;
    GActionMap *action_map;
    menu_layout_mode layout;
    bool show_arrow_indicators;
    widget_style_config style;
} menubar_config;
\end{lstlisting}

\subsection{Description des champs principaux}
\begin{center}
\begin{tabular}{|l|l|p{5cm}|}
\hline
\textbf{Champ} & \textbf{Type} & \textbf{Description} \\
\hline
kind & menu\_item\_kind & Type (Action, Sous-menu, Section) \\
\hline
children & menu\_item\_config* & Sous-éléments du menu \\
\hline
on\_activate & callback & Fonction appelée au clic \\
\hline
layout & menu\_layout\_mode & Horizontal ou Vertical \\
\hline
\end{tabular}
\end{center}

\bigskip

Dans la structure \texttt{menu\_item\_config}, le champ \texttt{label} définit le texte visible. Le champ \texttt{kind} est crucial : il détermine si l'élément est une action simple, un sous-menu ouvrant d'autres options, ou une section de séparation. 

Le champ \texttt{children} permet de lier d'autres structures \texttt{menu\_item\_config}, créant ainsi des menus en cascade. Le champ \texttt{enabled} permet d'activer ou de griser l'élément selon le contexte, tandis que \texttt{on\_activate} gère la fonction exécutée au clic.

Pour la \texttt{menubar\_config}, le champ \texttt{sections} organise la barre en blocs logiques (Fichier, Édition). Le champ \texttt{layout} définit si la barre est horizontale ou verticale, s'adaptant au design général. 

Le champ \texttt{show\_arrow\_indicators} ajoute un indice visuel pour les sous-menus, et le champ \texttt{style} assure la cohérence visuelle avec le reste de l'application via les classes CSS.

\subsection{Structure \texttt{menu\_section\_config}}
\begin{lstlisting}[language=C]
typedef struct {
    const char *label;
    const menu_item_config *items;
    size_t item_count;
} menu_section_config;
\end{lstlisting}
\begin{center}
\begin{tabular}{|l|l|>{\raggedright\arraybackslash}p{6cm}|}
\hline
\textbf{Champ} & \textbf{Type} & \textbf{Description} \\
\hline
label & const char* & Titre de la section (optionnel) \\
\hline
items & menu\_item\_config* & Tableau des éléments du menu pour cette section \\
\hline
item\_count & size\_t & Nombre d'éléments dans la section \\
\hline
\end{tabular}
\end{center}

\bigskip

Le champ \texttt{label} permet de définir un titre textuel pour un groupe d'éléments au sein d'un menu. Cela est particulièrement utile pour organiser des menus denses en séparant visuellement les catégories d'actions (par exemple, séparer "Édition" de "Formatage").

Le champ \texttt{items} est un pointeur vers un tableau de structures \texttt{menu\_item\_config}. Cette approche modulaire permet de construire chaque section indépendamment avant de les assembler dans la barre de menus globale, facilitant ainsi la maintenance du code.

Le champ \texttt{item\_count} indique au wrapper le nombre exact d'entrées contenues dans la section. Cela permet une allocation mémoire précise et assure que tous les éléments définis sont correctement rendus dans l'interface graphique finale.